\chapter{Project 5: Reasoning RL and Alignment (CS336 Assignment 5)}\label{ch:proj5}

The last CS336 assignment moves from pretraining to post-training. In the required part you
take a small but heavily pretrained base model, measure how well prompting alone solves
grade-school math word problems, and then improve it with on-policy policy gradients: first
standard GRPO, then a family of variants (constant normalisation, Dr.\ GRPO, rejection
fine-tuning, MaxRL), then off-policy training with importance reweighting and clipping, and
finally an estimator of your own design. An optional supplement covers the earlier half of the
post-training story: zero-shot evaluation of an 8B base model on knowledge, reasoning, chat and
safety benchmarks, supervised instruction tuning, red-teaming, and DPO on human preference data.

As with the other project guides (see the note at the start of \cref{ch:proj1}), this chapter
does not solve the assignment. It explains what each part asks in our own words, how to set up
the environment and the two-GPU system, which chapters teach what you need, and how to design,
debug and report the experiments. The mathematics you need is in \cref{ch:pg}, worked on
different examples from the handout's; the system design is in \cref{ch:rlsys}; the
evaluation machinery is in \cref{ch:eval}.

\begin{figure}[t]
\centering
\resizebox{\linewidth}{!}{%
\begin{tikzpicture}[st/.style={draw,rounded corners=3pt,minimum height=1.0cm,align=center,
  font=\footnotesize,text width=2.35cm},>={Stealth[length=2mm]}]
\node[st,fill=fdata!15] (p) at (0,0) {\textbf{1} prompting\\baselines (vLLM)};
\node[st,fill=fcompute!15] (g) at (2.9,0) {\textbf{2} derive \&\\build GRPO};
\node[st,fill=fcompute!15] (e) at (5.8,0) {\textbf{3} GRPO runs:\\seeds, LR, prompts};
\node[st,fill=fcomm!15] (v) at (8.7,0) {\textbf{4} on-policy\\variants};
\node[st,fill=fcomm!15] (o) at (11.6,0) {\textbf{5} off-policy:\\IS, clip, GSPO};
\node[st,fill=labpurple!15] (y) at (14.5,0) {\textbf{6} your own\\estimator};
\foreach \a/\b in {p/g,g/e,e/v,v/o,o/y} \draw[->,thick] (\a)--(\b);
\node[font=\small\bfseries,anchor=west] at (-1.2,-1.25) {optional supplement};
\node[st,fill=black!5] (z) at (1.5,-2.2) {zero-shot evals\\(8B base)};
\node[st,fill=fdata!15] (s) at (4.4,-2.2) {packed SFT\\+ evals};
\node[st,fill=warnred!12] (r) at (7.3,-2.2) {red-teaming};
\node[st,fill=fcomm!15] (d) at (10.2,-2.2) {DPO on HH\\+ evals};
\foreach \a/\b in {z/s,s/r,r/d} \draw[->,thick] (\a)--(\b);
\end{tikzpicture}}
\caption{The shape of the project. The required part (top) climbs from inference-only baselines
to a family of RL estimators; every stage after the second reuses the same training script with
different flags. The optional supplement (bottom) is a classic SFT-then-preference pipeline
with the same four evaluations after each stage.}
\label{fig:proj5:roadmap}
\end{figure}

% =====================================================================
\section{What the assignment asks, and why}\label{sec:proj5:overview}
\covers{CS336 A5 §1--2; CS336 L15; CS336 L16}

The testbed is deliberately small. The model is \textbf{OLMo-2-0425-1B}, a 1B-class base
model trained on roughly four trillion tokens, thousands of tokens per parameter, so it is far
more capable than its size suggests. The task is \textbf{GSM8K}, grade-school word problems
whose final answers are numbers that a verifier can check. Rewards are binary (correct or not),
with a separate format signal used only for logging. Responses are sampled with temperature
1.0, no nucleus truncation, and a 512-token cap, and generation stops at the closing answer tag.
This is RLVR in its cleanest form (\cref{ch:grpo}): no reward model, no human labels, just a
grader.

The learning goals are three. (1) Understand \emph{which objective} each estimator optimises,
not just how to code it: the assignment asks you to derive variances, bias factors and the
prompt reweightings induced by different normalisers (the method is in
\cref{sec:pg:baselines,sec:pg:normalisation}). (2) Build a correct and efficient RL training
step around a Hugging Face model and a vLLM rollout server on separate GPUs. (3) Do RL
\emph{science} under noise: multiple seeds, controlled ablations, honest confidence statements.
The handout stresses that RL runs vary greatly across seeds, and the grading reflects that.

\begin{table}[ht]
\centering\small
\begin{tabularx}{\linewidth}{@{}lX@{}}
\toprule
Part & What you produce\\
\midrule
Prompting & an evaluation script; a breakdown of outputs by (format, answer) reward; qualitative behaviour of three prompt styles\\
On-policy GRPO & derivation answers; tokenisation, log-prob/entropy, reward, advantage, loss and aggregation functions; a train step with accumulation; a full training script; 4-seed runs and plots; learning-rate and prompt sweeps\\
Variants & normaliser and aggregation options; a train step that skips zero-advantage rows; 4-seed comparison of four variants\\
Off-policy & token- and sequence-level reweighted losses with clipping; a 32$\times$ off-policy schedule; 4-seed comparison with clip-fraction logging\\
Your estimator & one change to an existing estimator, motivated, run with seeds\\
Supplement & MMLU/GSM8K/AlpacaEval/SimpleSafetyTests evaluation scripts; packed SFT data loader and trainer; red-team report; HH loader; DPO loss and trainer; before/after evaluations\\
\bottomrule
\end{tabularx}
\caption{Deliverables at a glance, paraphrased. The handout is authoritative on exact
requirements and point values.}
\label{tab:proj5:deliverables}
\end{table}

% =====================================================================
\section{Setting up}\label{sec:proj5:setup}
\covers{CS336 A5 §1, §3.1, §4.2.1--4.2.2}

\subsection{Repository and environment}

Clone the public assignment repository. Its layout matters because the tests drive your code
through it:
\begin{itemize}
\item \texttt{cs336\_alignment/}: your package. It ships with starter utilities only: a
  checkpoint loader for the Hugging Face model, a vLLM server wrapper that can also receive
  weights over NCCL, the math grader used for rewards, prompt templates, and Modal helpers.
  Everything else you write.
\item \texttt{tests/}: the required RL tests and the optional supplement tests. They call
  thin hook functions in \texttt{tests/adapters.py}; your job is to make each hook call your
  implementation. Read the adapters and the test files first, since they define exactly which tensor
  shapes, dtypes and edge cases are checked.
\item \texttt{data/}: GSM8K train/test splits and the supplement's benchmarks (MMLU, AlpacaEval
  references, SimpleSafetyTests, Anthropic HH).
\item \texttt{scripts/}: the AlpacaEval annotator configuration and a safety-judging script for
  the supplement.
\end{itemize}
Dependencies are managed with \texttt{uv}, pinned to Python 3.12. FlashAttention does not
install cleanly in one pass, so the README installs everything else first and then the full set.
The GPU extras pin a specific vLLM release, so do not upgrade vLLM independently: its weight-sync
and generation APIs change between versions.

\begin{lstlisting}[style=shell,caption={Environment bring-up and a first test run.},label={lst:proj5:setup}]
git clone https://github.com/stanford-cs336/assignment5-alignment
cd assignment5-alignment
uv sync --no-install-package flash-attn && uv sync   # two steps
uv run pytest tests/test_grpo.py         # initially: NotImplementedError
python BOOK/code/projects/a5/env_check.py --need-gpus 2 --min-gpu-gb 40
\end{lstlisting}

The environment checker in \texttt{code/projects/a5} (part of this book, not the course)
reports the Python version, whether torch, transformers, vLLM, FlashAttention, wandb and
math-verify import, the GPU inventory with bf16 support, free disk, and whether a W\&B key is
set. It touches nothing else. Run it on the machine that will do the work, before spending GPU
hours.

\subsection{Compute: Modal, B200s, and your own GPUs}

The course runs jobs on Modal. A helper module builds a container image from the lockfile,
requests two B200 GPUs per job, mounts your code and data, and maps a list of command lines
across containers, so four seeds become four concurrent containers. You set your own identifier
in that module and store a W\&B API key as a Modal secret. The handout budgets GPU hours per
experiment problem (\cref{fig:proj5:budget}); plan your sweeps against those numbers, not
against your enthusiasm.

\begin{figure}[ht]
\centering
\begin{tikzpicture}
\begin{axis}[xbar,width=0.8\linewidth,height=5.4cm,bar width=7pt,xmin=0,xmax=10,
  xlabel={B200 GPU-hours budgeted by the handout},
  symbolic y coords={DPO: supplement,SFT: supplement,off-policy variants,on-policy variants,prompt ablation,learning-rate sweep,GRPO 4 seeds},
  ytick=data,nodes near coords,nodes near coords style={font=\scriptsize},
  tick label style={font=\scriptsize},label style={font=\small}]
\addplot[fill=sysblue!55] coordinates {(2,GRPO 4 seeds) (4,learning-rate sweep) (4,prompt ablation) (8,on-policy variants) (8,off-policy variants) (3,SFT: supplement) (1,DPO: supplement)};
\end{axis}
\end{tikzpicture}
\caption{GPU-hour budgets stated in the handout for the experiment problems: 26 B200-hours for
the required part and 4 for the supplement. Budget for debugging runs on top: most of the
risk is in getting a correct script, not in the final sweeps.}
\label{fig:proj5:budget}
\end{figure}

\begin{napkinmath}[translating B200 hours to your hardware]
A B200 offers roughly twice the dense bf16 tensor throughput of an H100, and more HBM bandwidth
and capacity; treat one B200-hour as about two H100-hours for training-bound work. Rollouts are
memory-bound (\cref{ch:inference}), so for them the ratio tracks HBM bandwidth instead. The
required experiments then cost on the order of 50 H100-hours, or a few hundred hours on a
24\,GB consumer card if the model and engine fit at all. Before you start, measure your own
step time on 5 steps and multiply out the planned runs. The \texttt{step\_time} breakdown of
\cref{tab:pg:metrics} tells you which phase to optimise first.
\end{napkinmath}

\subsection{The two-GPU system}

The recommended layout is exactly \cref{fig:pg:system}: the Hugging Face policy and its AdamW
state on GPU~0, a vLLM server on GPU~1, and an NCCL weight-transfer group between them. Each
training step syncs the policy's weights into vLLM (which also resets caches that depended on
the old weights), generates rollouts with the stop string included in the output, scores them
on the CPU, and runs the micro-batched update on GPU~0. \Cref{fig:proj5:timeline} shows the
resulting timeline. In the synchronous design, each GPU idles while the other works. That is
acceptable for this project; \cref{ch:rlsys} shows how asynchronous schedules recover the idle
time and what they cost in staleness.

\begin{figure}[t]
\centering
\begin{tikzpicture}[x=0.66cm,y=1.05cm]
\node[anchor=east,font=\scriptsize] at (-0.1,1) {GPU 1 (vLLM)};
\node[anchor=east,font=\scriptsize] at (-0.1,0) {GPU 0 (trainer)};
\node[anchor=east,font=\scriptsize] at (-0.1,-1) {CPU};
\foreach \s in {0,1} {
  \pgfmathsetmacro{\o}{\s*9.4}
  \fill[fmemory!55] (\o+0.2,0.72) rectangle (\o+4.4,1.28); \node[font=\scriptsize] at (\o+2.3,1) {generate rollouts};
  \fill[napkin!45] (\o+4.4,-1.28) rectangle (\o+5.4,-0.72); \node[font=\scriptsize] at (\o+4.9,-1) {grade};
  \fill[fcompute!45] (\o+5.4,-0.28) rectangle (\o+8.8,0.28); \node[font=\scriptsize] at (\o+7.1,0) {train step};
  \fill[warnred!55] (\o+8.8,-0.28) rectangle (\o+9.4,1.28); \node[font=\tiny,rotate=90] at (\o+9.1,0.5) {sync};
}
\draw[->] (0,-1.6)--(19.2,-1.6) node[below left,font=\scriptsize]{time};
\node[font=\scriptsize,text=black!50] at (2.3,0) {idle};
\node[font=\scriptsize,text=black!50] at (7.1,1) {idle};
\end{tikzpicture}
\caption{Synchronous on-policy RL on two GPUs. Generation, grading, training and weight sync
alternate; each GPU idles for roughly the other's busy time. With 32$\times$ off-policy reuse
(Part~5), the trainer segment grows and generation becomes rarer.}
\label{fig:proj5:timeline}
\end{figure}

\begin{pitfall}[things that break the two-GPU setup]
(i) The vLLM server and the trainer must load the \emph{same} weights and tokenizer revision.
(ii) Weight sync must happen \emph{before} each generation, including the first one if you
changed anything at load time. Check it with the log-probability test of the debugging list in
\cref{sec:pg:loop}. (iii) Generation stop strings must be included in the returned text, or the
grader will fail to find the closing tag and every reward will be zero. (iv) Give vLLM an
explicit GPU memory fraction. On a shared GPU it will otherwise claim most of the memory for its
KV pool.
\end{pitfall}

\subsection{What the provided utilities do}\label{sec:proj5:utils}

Four starter modules carry the infrastructure, so you can focus on the algorithm. Learn their
contracts before writing anything.

\begin{description}
\item[Model loading.] A helper loads the policy and tokenizer with Transformers in bf16, using
  FlashAttention-2 on GPU and eager attention on CPU (useful for unit tests). Trained models are
  written back with \code{save\_pretrained} for both model and tokenizer, so a checkpoint
  directory is self-contained and loadable by vLLM.
\item[The vLLM server wrapper.] It starts a vLLM HTTP server as a subprocess pinned to one GPU,
  with a seed and a GPU-memory fraction, and generates completions for a list of prompt strings
  and a dictionary of sampling parameters. Each completion carries text, token IDs and the finish
  reason. Two more calls create an NCCL group between the trainer's device and the server and
  push the current policy weights into it. Because the server returns token IDs, you can follow
  \cref{sec:pg:batch} and train on exactly what was sampled.
\item[The grader.] Reward functions take a response and a ground-truth answer and return a
  dictionary with a total reward, a format reward and an answer reward. There is one for the
  tagged prompts (it looks inside the answer tags) and one for the boxed-answer prompt. Ground
  truths in GSM8K end with a marker followed by the number; split on it and strip.
\item[Modal helpers.] A small app definition builds the image from the lockfile, requests two
  GPUs, mounts code and data, and maps a list of command lines over containers. Your scripts stay
  ordinary Python with command-line flags; the helper only launches them.
\end{description}

\begin{systemslens}[the server boundary is an API boundary]
Running vLLM as a separate server process, rather than inside the trainer, isolates its CUDA
context, memory pool and kernels from yours, at the cost of an HTTP hop per generation call and
an explicit NCCL group for weights. This is the same separation that production RL frameworks
make between the rollout and training worker groups (\cref{sec:rlsys:placement}); the
assignment gives you the smallest working instance of it. Batch your generation calls: one
request with 256 prompts lets vLLM's continuous batcher (\cref{ch:serving}) keep the GPU busy,
while 256 sequential requests leave it idle between them.
\end{systemslens}

\subsection{Structuring your code}\label{sec:proj5:structure}

The same functions serve the tests, the evaluation script, five experiment families and the
supplement, so organise them for reuse from the start (\cref{fig:proj5:modules}). Keep pure
tensor functions (masks, log-probs, advantages, losses, aggregation) free of I/O so they are
trivially unit-testable on CPU; keep anything that touches GPUs, servers or files in thin
drivers. Expose every algorithmic switch the handout mentions (baseline, normaliser,
aggregation, reweighting method, clip range, number of updates per rollout batch) as a
command-line flag of one training script. The ablations then become a list of commands rather
than a set of forked scripts.

\begin{figure}[ht]
\centering
\begin{tikzpicture}[m/.style={draw,rounded corners=2pt,font=\scriptsize,align=center,minimum height=7mm,text width=24mm},
  >={Stealth[length=1.6mm]}]
\node[m,fill=fcompute!12] (core) at (0,0) {\textbf{pure functions}\\masks, log-probs,\\advantages, losses};
\node[m,fill=fdata!12] (data) at (-3.4,0) {\textbf{data \& prompts}\\loaders, templates};
\node[m,fill=napkin!15] (grade) at (3.4,0) {\textbf{grading}\\provided reward fns};
\node[m,fill=fmemory!20] (rl) at (-1.7,-1.6) {\textbf{RL driver}\\server, sync, loop,\\logging, ckpt};
\node[m,fill=fmemory!20] (ev) at (1.7,-1.6) {\textbf{eval driver}\\generate, parse,\\serialise};
\node[m,fill=black!5] (tests) at (0,1.5) {\textbf{tests}\\adapters $\to$ pure fns};
\node[m,fill=fcomm!12] (sup) at (5.1,-1.6) {\textbf{supplement}\\SFT, DPO drivers};
\node[m,fill=labpurple!12] (launch) at (-5.1,-1.6) {\textbf{launch}\\flag grid $\to$ Modal};
\draw[->] (tests)--(core); \draw[->] (rl)--(core); \draw[->] (ev)--(core);
\draw[->] (rl)--(data); \draw[->] (ev)--(grade); \draw[->] (rl)--(grade); \draw[->] (sup)--(ev);
\draw[->] (launch)--(rl);
\end{tikzpicture}
\caption{A layout that keeps the tested core small and pure. Drivers depend on the core,
never the other way round. The evaluation driver is shared by the prompting baselines, RL
validation and the supplement.}
\label{fig:proj5:modules}
\end{figure}

% =====================================================================
\section{Problem-by-problem map}\label{sec:proj5:map}

\Cref{tab:proj5:map} lists each problem, paraphrased, with the parts of the book that teach
what it needs. Where it says ``method'', the referenced section develops the technique on a
different example; apply it to the handout's setting yourself.

{\small
\begin{longtable}{@{}>{\raggedright\arraybackslash}p{0.27\linewidth}>{\raggedright\arraybackslash}p{0.38\linewidth}>{\raggedright\arraybackslash}p{0.28\linewidth}@{}}
\caption{CS336 A5 problems and where the book covers the background.}\label{tab:proj5:map}\\
\toprule
\textbf{Problem} & \textbf{What it asks (paraphrased)} & \textbf{Read}\\
\midrule\endfirsthead
\toprule
\textbf{Problem} & \textbf{What it asks (paraphrased)} & \textbf{Read}\\
\midrule\endhead
{\raggedright\texttt{\footnotesize prompting\_\allowbreak baselines}} & evaluate three prompt styles with vLLM; bucket outputs by format and answer reward; inspect parse failures & \cref{sec:eval:verifiers,sec:eval:extract,sec:eval:run}; \cref{sec:scaling:cot}\\
{\raggedright\texttt{\footnotesize baseline\_\allowbreak calcs}} & variance of a score-function estimator with and without a baseline on a toy policy & method: \cref{sec:pg:baselines}\\
{\raggedright\texttt{\footnotesize tokenize\_\allowbreak prompt\_\allowbreak and\_\allowbreak output}} & prompt+response layout with shifted labels and a response mask & \cref{sec:pg:batch}, \cref{fig:pg:layout}\\
{\raggedright\texttt{\footnotesize get\_\allowbreak response\_\allowbreak log\_\allowbreak probs}} & per-token log-probs and optional entropy from a causal LM & \cref{sec:pg:logprobs}\\
{\raggedright\texttt{\footnotesize compute\_\allowbreak rollout\_\allowbreak rewards}} & score responses with the provided grader and log statistics & \cref{sec:eval:grade}; \cref{ch:grpo}\\
{\raggedright\texttt{\footnotesize compute\_\allowbreak group\_\allowbreak normalized\_\allowbreak rewards\_\allowbreak *}} & group baseline with std, none or mean normalisers & \cref{sec:pg:practical,sec:pg:reweight}\\
{\raggedright\texttt{\footnotesize compute\_\allowbreak policy\_\allowbreak gradient\_\allowbreak loss\_\allowbreak *}} & on-policy, token-reweighted, clipped and sequence-level (GSPO) per-token losses & \cref{sec:pg:reinforce,sec:pg:offpolicy}; \cref{sec:grpo2:clip,sec:grpo2:more}\\
{\raggedright\texttt{\footnotesize aggregate\_\allowbreak loss\_\allowbreak *}} & sequence-mean and constant aggregation & \cref{sec:pg:length,sec:pg:aggregation}\\
{\raggedright\texttt{\footnotesize grpo\_\allowbreak train\_\allowbreak step\_\allowbreak *}} & micro-batched step with clipping, metrics, zero-advantage pruning, off-policy arguments & \cref{prop:pg:accum}; \cref{sec:pg:loop}\\
{\raggedright\texttt{\footnotesize grpo\_\allowbreak experiments\_\allowbreak *}} & full training script; 4 seeds; learning-rate and prompt sweeps; variant and off-policy comparisons & \cref{sec:pg:loop}; \cref{ch:grpo2}; \cref{sec:proj5:experiments}\\
{\raggedright\texttt{\footnotesize think\_\allowbreak about\_\allowbreak *}} & short discussions of length normalisation, RFT, advantage normalisation, importance reweighting & \cref{sec:pg:length,sec:pg:rft,sec:pg:reweight,sec:pg:is}\\
{\raggedright\texttt{\footnotesize derive\_\allowbreak *}} & reweightings induced by normalisers; surrogate objective of a grouped token ratio & method: \cref{sec:pg:reweight,sec:pg:is}\\
{\raggedright\texttt{\footnotesize try\_\allowbreak your\_\allowbreak own}} & propose and test one change to an estimator & \cref{sec:proj5:own}\\
\midrule
\multicolumn{3}{@{}l}{\emph{Supplement}}\\
{\raggedright\texttt{\footnotesize *\_\allowbreak baseline}}, {\raggedright\texttt{\footnotesize *\_\allowbreak sft}} evals & MMLU letter parsing, GSM8K last-number parsing, AlpacaEval win rates, safety judging & \cref{sec:eval:mcq,sec:eval:judge}; \cref{ch:alignment}\\
{\raggedright\texttt{\footnotesize look\_\allowbreak at\_\allowbreak sft}}, {\raggedright\texttt{\footnotesize data\_\allowbreak loading}} & inspect instruction data; packed fixed-length SFT dataset and batch iterator & \cref{sec:instruct:packing}; \cref{ch:alignment}\\
{\raggedright\texttt{\footnotesize sft\_\allowbreak script}}, {\raggedright\texttt{\footnotesize sft}} & full-parameter SFT of an 8B model with accumulation and a cosine schedule & \cref{sec:instruct:train}; \cref{ch:optim}; \cref{sec:proj5:supplement}\\
{\raggedright\texttt{\footnotesize red\_\allowbreak teaming}} & probe for misuse and report methodology & \cref{ch:alignment}\\
{\raggedright\texttt{\footnotesize look\_\allowbreak at\_\allowbreak hh}}, {\raggedright\texttt{\footnotesize dpo\_\allowbreak loss}}, {\raggedright\texttt{\footnotesize dpo\_\allowbreak training}} & HH preference loader; per-instance DPO loss; two-GPU DPO training with validation accuracy & \cref{ch:dpo}; \cref{sec:proj5:supplement}\\
\bottomrule
\end{longtable}}

% =====================================================================
\section{Guidance for the prompting baselines}\label{sec:proj5:prompting}
\covers{CS336 A5 §3}

This part is short, but it calibrates everything after it. Three prompt styles are compared:
the bare question with an instruction to box the answer, an R1-Zero-style prompt that opens a
thinking tag and expects an answer tag, and a three-shot version of the latter. Each has a
matching grader in the provided module, and the boxed-answer style uses a different parser.
Use the stop string only for the tagged prompts.

\begin{itemize}
\item Write one evaluation function that takes a prompt template, a grader and sampling
  parameters, and serialises \texttt{(prompt, response, rewards, finish\_reason)} for every
  example. You will reuse it inside the RL loop for validation.
\item Count outcomes in the three (format, answer) buckets, then \emph{read} the
  format-correct-but-wrong and format-wrong buckets. The handout asks how many are actually
  correct but unparsed. That number estimates your reward noise, the floor on what RL can fix.
\item Note what the base model does besides answering: continuing the conversation,
  generating new questions, emitting code. These behaviours are what RL will suppress or
  amplify, and the prompt determines which ones appear.
\item Record throughput (examples per second) and the response-length distribution. Together
  they set your rollout budget for every later part.
\end{itemize}

% =====================================================================
\section{Guidance for building GRPO}\label{sec:proj5:grpo}
\covers{CS336 A5 §4.1--4.2}

\paragraph{Derivations.} The written problems test the ideas of \cref{sec:pg:baselines}. For
the variance problem, write the estimator's first and second moments for the given policy and
reward, exactly as \cref{fig:pg:variance} does for a different bandit, and check your algebra
numerically: sample $10^5$ actions and compare with your formula, as Lab~24 does. For the
reweighting derivations, follow the large-group argument of \cref{sec:pg:reweight}: identify
what each group statistic converges to as $G\to\infty$, factor it out of the expectation, and
read off the weight.

\paragraph{Components.} Write each function so that it is testable in isolation, then make the
adapters thin. Invariants worth asserting in your own tests:
\begin{itemize}
\item tokenisation: \code{response\_mask.sum(1)} equals the response token counts; the first
  masked label equals the first response token; nothing is inserted between prompt and response
  (\cref{fig:pg:layout});
\item log-probs: every value is $\le0$; entropies lie in $[0,\log V]$; outputs are float32 even
  when the model runs in bf16 (\cref{lst:pg:logprobs});
\item advantages: zero for groups with identical rewards; each group's advantages sum to zero
  under a mean baseline; finite everywhere, thanks to $\epsilon$;
\item losses: with all ratios equal to one, the reweighted and clipped losses equal the
  on-policy loss; the clip fraction is zero.
\end{itemize}

\paragraph{The train step.} The step must give the same gradient whatever the number of
accumulation steps. That is \cref{prop:pg:accum}: every micro-batch must be normalised by
quantities of the \emph{whole} rollout batch. Test it directly. Compute the gradient with one
micro-batch and with eight on a tiny batch, and compare. Clip the gradient norm once, after the
last micro-batch, and log its pre-clip value.

\paragraph{Memory.} OLMo-2-1B in bf16 with fp32 AdamW state needs about 16 bytes per
parameter, around 24\,GB, before activations, so the trainer GPU is not tight. The logits are
the largest transient tensor: work through the napkin math in \cref{sec:pg:logprobs} for your
micro-batch size and vocabulary before choosing it. If you run out of memory, shrink the
micro-batch and keep the accumulation correct, rather than shortening responses, which changes
the task.

% =====================================================================
\section{Designing the experiments}\label{sec:proj5:experiments}
\covers{CS336 A5 §4.3, §5.4, §6.4}

\paragraph{Correctness first.} Before any sweep, run about 50 steps on one seed and confirm
three things: validation reward rises (it may start near zero and wait a few steps for the
first successes), entropy declines gradually rather than collapsing, and sampled rollouts read
sensibly. The handout's suggested hyperparameters are tuned to make this work on most seeds; if
yours does not learn with them, suspect a bug before suspecting the hyperparameters, and work
down the checklist of \cref{sec:pg:loop}.

\paragraph{Seeds and plots.} The handout asks for four seeds per configuration and plots that
show the spread. Name runs so that seeds group together, log JSON lines per step, and aggregate
with intervals. \code{code/projects/a5/run\_matrix.py} expands a settings$\times$seeds grid into
command lines (one per Modal container), and \code{seeds\_ci.py} turns per-seed logs into
mean, interval and min/max per step (\cref{fig:proj5:matrix}). Plot the individual runs too:
with $n=4$, a normal-approximation interval is optimistic.

\begin{figure}[t]
\centering
\begin{tikzpicture}[c/.style={draw,minimum width=11mm,minimum height=6mm,font=\scriptsize}]
\node[font=\small\bfseries] at (2.6,1.0) {seed};
\foreach \s [count=\i from 0] in {0,1,2,3} \node[font=\scriptsize] at (\i*1.2+0.8,0.5) {\s};
\foreach \v [count=\j from 0] in {GRPO (seq),GRPO (const),Dr.\ GRPO,RFT,MaxRL} {
  \node[anchor=east,font=\scriptsize] at (0.1,-\j*0.7) {\v};
  \foreach \i in {0,1,2,3} \node[c,fill=fcompute!12] at (\i*1.2+0.8,-\j*0.7) {run};}
\draw[->,thick] (5.4,-1.4)--(6.6,-1.4);
\node[draw,rounded corners=2pt,fill=keygreen!12,align=center,font=\scriptsize,text width=36mm] at (8.6,-1.4)
  {\code{seeds\_ci.py}: per step, mean $\pm1.96\,s/\sqrt{n}$, min, max; one curve and band per row};
\end{tikzpicture}
\caption{An ablation is a grid of (variant $\times$ seed) runs that differ in one setting.
Group by variant, aggregate across seeds, and compare bands, not single curves.}
\label{fig:proj5:matrix}
\end{figure}

\paragraph{Controlled comparisons.} Each comparison should change one thing. The handout tunes
the learning rate on standard GRPO and reuses it for the variants. As \cref{sec:pg:loop} notes,
that makes a variant's win convincing but its loss inconclusive. Say so in the write-up. For the
variant sweep, implement zero-advantage pruning and verify that it leaves the gradient unchanged
(\cref{sec:pg:aggregation}); for rejection fine-tuning it removes most of the training compute.
For the off-policy sweep, the training batch shrinks to 1/32 of the rollout batch, so each
rollout batch feeds 32 optimizer steps. Compute the old log-probabilities once, before the first
of them, and log the clip fraction and approximate KL for every step. The off-policy runs also
change the wall-clock balance of \cref{fig:proj5:timeline}: measure and report it, because a method
that is better per rollout but slower per hour may still lose.

\paragraph{What to log.} At minimum: loss, pre-clip gradient norm, token entropy, train reward
and format reward, validation reward and format reward, validation response length, plus clip
fraction and zero-advantage fraction where relevant (\cref{tab:pg:metrics}). Every 40 steps or
so, dump a few rollouts to disk and read them. Reward hacking and format drift are visible in
text long before they show up in a curve.

\subsection{Failure signatures}\label{sec:proj5:failures}

RL bugs rarely crash; they show up as odd curves. \Cref{tab:proj5:failures} lists common
signatures and the first thing to check. Most of them map to an invariant from
\cref{sec:proj5:grpo} or to the checklist in \cref{sec:pg:loop}.

{\small
\begin{longtable}{@{}>{\raggedright\arraybackslash}p{0.3\linewidth}>{\raggedright\arraybackslash}p{0.3\linewidth}>{\raggedright\arraybackslash}p{0.33\linewidth}@{}}
\caption{Symptoms, likely causes and first checks for RL runs.}\label{tab:proj5:failures}\\
\toprule
\textbf{Symptom} & \textbf{Likely cause} & \textbf{First check}\\
\midrule\endfirsthead
\toprule
\textbf{Symptom} & \textbf{Likely cause} & \textbf{First check}\\
\midrule\endhead
train reward exactly 0 for many steps & grader never fires; stop string stripped; wrong ground-truth split & grade ten responses by hand\\
reward rises, validation flat & validation uses a different template, parser or sampling settings & diff the two code paths\\
loss decreases but nothing improves & sign error (ascending instead of descending) or mask misaligned & flip test: negate advantages, the reward should fall\\
reward rises then collapses & learning rate too high; off-policy updates drift & gradient norm and clip fraction around the collapse\\
entropy drops to near 0 within a few steps & too-large steps; very small groups & lower LR; inspect rollouts for repetition\\
responses grow to the length cap & length-biased aggregation; truncated responses scored as failures & share of truncated responses; aggregation mode\\
results differ with gradient-accumulation steps & micro-batches normalised locally & 1-vs-$k$ gradient comparison (\cref{prop:pg:accum})\\
model never changes in vLLM & weight sync not called or not applied & log-prob of a fixed string before and after sync\\
GPU memory spikes on long batches & full-vocabulary fp32 logits & chunking and response-only slicing (\cref{sec:pg:logprobs})\\
off-policy runs identical to on-policy & old log-probs recomputed with current weights after each update & compute old log-probs once, before the first update\\
\bottomrule
\end{longtable}}

\subsection{Your own estimator}\label{sec:proj5:own}
\covers{CS336 A5 §7}

The final problem rewards a clean experiment more than a clever idea. Choose a change that the
chapter's framework can motivate: a different baseline (for example leave-one-out with a
different group statistic), a different prompt reweighting (derive its target with
\cref{sec:pg:reweight} before running it), a different treatment of importance weights
(truncation, masking, or a different geometric aggregation), or a different token weighting.
Write the predicted effect in one sentence \emph{before} running it, change exactly one thing
relative to a variant you already ran, use the same seeds and budget, and report a result even
if it is negative.

% =====================================================================
\section{The supplement: SFT, evaluation, red-teaming and DPO}\label{sec:proj5:supplement}
\covers{CS336 A5 supplement §2--6; CS336 L15}

The supplement uses an 8B base model, Llama~3.1~8B, with a 70B instruction-tuned model as an
automatic judge. Weights and the instruction-tuning data live on a shared Modal volume; write
your own checkpoints only to your private results volume. Everything is built from scratch except
inference (vLLM) and model loading (Transformers); no trainer classes.

\paragraph{Evaluations.} Four benchmarks probe four abilities: MMLU (knowledge; parse a letter
from a fixed-format sentence), GSM8K (reasoning; take the last number in the output), AlpacaEval
(chat quality; win rate and length-controlled win rate against a reference model, judged by the
70B model), and SimpleSafetyTests (safety; the share of responses a judge deems safe). All use
greedy decoding. Build one harness that formats, generates, parses and serialises, and plug in a
parser per task. \Cref{ch:eval} gives the MCQ and judge machinery, and \cref{ch:alignment} builds
a multi-benchmark harness with a mock judge you can test offline. The prompt template changes
after SFT (the training template replaces the zero-shot system prompt), so keep templates in
files, not in code.

\paragraph{Packed SFT.} The data are single-turn prompt/response pairs. Format each with the
training template, join them into one token stream separated by the end-of-text token, and cut
it into non-overlapping fixed-length windows. That gives no padding, maximal throughput, and
documents spanning window boundaries (\cref{sec:instruct:packing} contrasts this with
per-example padding and loss masking). Then write a trainer with gradient accumulation, a
warmup-plus-cosine schedule, weight decay and clipping (\cref{ch:optim,ch:pretrain}).

\begin{napkinmath}[full fine-tuning an 8B model on one GPU]
Weights in bf16 take 16\,GB, gradients in bf16 16\,GB, and AdamW moments in fp32 64\,GB, for
about 96\,GB before activations, or 128\,GB with fp32 master weights. That fits on one
high-memory accelerator (B200-class, roughly 180--192\,GB) at sequence length 512 with a small
micro-batch, but not on an 80\,GB H100. There you would need ZeRO/FSDP across GPUs
(\cref{ch:sharding}), an 8-bit optimizer, or LoRA (\cref{ch:instruct}). The DPO part holds a
trainable policy plus a frozen reference: putting the reference on its own GPU and using RMSprop,
which keeps one state tensor instead of two, as the handout suggests, keeps the policy GPU
around 64--96\,GB.
\end{napkinmath}

\paragraph{Red-teaming.} Treat it as an evaluation methodology, not a stunt. Pick concrete misuse
categories, record each attempt (strategy, prompt family, outcome, time spent), and summarise
what worked. Keep outputs inside your private results and report patterns, not harmful
completions.

\paragraph{DPO.} Load the four HH training subsets, drop conversations with more than one human
turn, split each into instruction plus chosen and rejected responses, and remember each
example's source subset. The per-instance loss needs four sequence log-probabilities: policy and
reference, chosen and rejected. The prompt's contribution cancels in each difference, so you can
score full concatenated strings. \Cref{ch:dpo} derives the objective, shows how to cache reference
log-probabilities, and explains the implicit-reward accuracy you should track on a held-out set.
Evaluate the DPO model on all four benchmarks, and look for an ``alignment tax'' on GSM8K and MMLU.

% =====================================================================
\section{The write-up}\label{sec:proj5:writeup}

The written deliverables are where this assignment is graded on judgement, so plan them like a
short paper.
\begin{itemize}
\item \textbf{One figure per question.} Each comparison gets a figure with one band per variant,
  plus a panel or appendix plot of individual seeds. Put the metric's definition, the number
  of seeds and the interval type in the caption.
\item \textbf{Claims sized to evidence.} With four seeds, ``A beats B'' needs bands that do not
  overlap over the last tenth of training; otherwise write ``indistinguishable at this budget''.
  State tuning asymmetries explicitly.
\item \textbf{Mechanism, not just outcome.} When a variant wins or loses, connect it to a logged
  quantity: entropy, response length, clip fraction, zero-advantage fraction. Then connect that
  quantity to the estimator's math (which prompts or tokens it upweights, \cref{sec:pg:reweight}).
\item \textbf{Systems numbers.} Report the step-time breakdown and wall-clock per configuration.
  Off-policy reuse and pruning are as much systems interventions as algorithmic ones.
\item \textbf{Qualitative evidence.} Quote two or three rollouts before and after training for the
  headline run, chosen to show a pattern (shorter or longer reasoning, format adherence), not a
  cherry-picked success.
\end{itemize}

\subsection{The supplement in more detail}\label{sec:proj5:supdetail}

\paragraph{Inspecting data.} For both the SFT data and HH, read random samples, not the first
lines of the file. For SFT, note which classic tasks hide inside the conversational prompts
(question answering, rewriting, summarisation, classification, code) and how good the responses
are. For HH, compare chosen and rejected responses within each subset. The ``harmless'' subset
often prefers refusals or deflections, the ``helpful'' subsets prefer substance, and the two
pull in opposite directions in places. That tension is the subject of \cref{ch:alignment}.

\paragraph{Win rates.} AlpacaEval reports a raw win rate and a length-controlled win rate that
regresses out the judge's preference for longer answers. Report both, and when they disagree,
inspect response lengths: an SFT or DPO model that ``wins'' only on raw win rate has mostly
learned to write more. Judges have positional and stylistic biases too (\cref{sec:eval:judge}).
Spot-check annotations where you disagree, as the handout asks.

\paragraph{Safety judging.} A judge model labels responses safe or unsafe, which is a proxy.
Look at both error types: unsafe responses judged safe, and refusals of benign requests, which
the metric does not penalise but users do. Over-refusal after safety SFT or DPO is a real
regression.

\paragraph{DPO training details.} Start from the handout's suggested values for effective batch
size, $\beta$ and learning rate, with one pair per micro-batch and gradient accumulation. Track the
loss, the implicit-reward accuracy on a held-out set of about 200 pairs, and the reward margin.
Checkpoint on the best validation accuracy rather than the last step: DPO overfits small
preference sets quickly (\cref{ch:dpo} discusses why, including the tendency of both chosen
and rejected log-probabilities to fall). Validate your loss on a pair where you can compute the four
log-probabilities by hand from a tiny randomly initialised model before running it at 8B.

% =====================================================================
\section{Background: what open reasoning-RL recipes teach}\label{sec:proj5:cases}
\covers{CS336 L16}

The assignment isolates the estimator, but the recipes behind public reasoning models show what
matters around it. These lessons are worth keeping in mind when you interpret your curves.
\begin{itemize}
\item \textbf{The base model sets the ceiling.} DeepSeek-R1-Zero (DeepSeek-AI, 2025) ran GRPO
  with accuracy and format rewards directly on a strong base model, and reasoning behaviour
  emerged from RL alone. Follow-up analyses (Liu et al., 2025) found that much of the ``aha''
  behaviour was already present in the base model and that part of the length growth came from
  the objective's length bias (\cref{sec:pg:length}). The assignment's choice of a model trained
  on thousands of tokens per parameter follows the same logic: RL sharpens what pretraining put
  there.
\item \textbf{Data curation is half the algorithm.} Several reports filter prompts by difficulty,
  keeping problems the model sometimes but not always solves, which is the regime where group
  baselines carry signal (\cref{sec:pg:practical}). Kimi~k1.5 (Kimi Team, 2025) also used a
  curriculum and sampled problems in proportion to their failure rate. Qwen3's report describes
  reasoning RL on a few thousand carefully filtered problems.
\item \textbf{Length is a first-class quantity.} Kimi~k1.5 added an explicit length reward to
  compress reasoning; others cap or penalise overlong responses (\cref{ch:grpo2}). In your runs,
  response length is the first diagnostic to plot next to reward.
\item \textbf{Simple pipelines win.} The open recipes replaced process reward models and tree
  search with outcome rewards and plain policy gradients. Your job in this assignment is to
  understand that plain policy gradient deeply enough to modify it with intent.
\end{itemize}

% =====================================================================
\section{A practical plan}\label{sec:proj5:plan}

The fastest path through the project front-loads correctness and back-loads compute.
\begin{enumerate}
\item \textbf{Afternoon 1, CPU only.} Read the adapters and tests. Implement the pure functions and
  pass their tests. Write your own invariant tests. Reproduce the derivation results numerically on
  toy policies.
\item \textbf{Afternoon 2, one GPU pair.} Bring up vLLM; run the prompting evaluation on 100
  examples, then on the full test set. Measure throughput and length distributions.
\item \textbf{Day 3.} Implement the train step; run the 1-vs-$k$ gradient check and the
  log-probability match against vLLM. Run 50 steps of standard GRPO on one seed.
\item \textbf{Day 4.} Launch four seeds of standard GRPO; while they run, write the plotting and
  aggregation scripts. Start the learning-rate sweep once the headline curve looks right.
\item \textbf{Days 5--6.} Variant and off-policy sweeps from one script with flags; your own
  estimator, designed while the sweeps run.
\item \textbf{Day 7.} Write-up. Re-run anything whose bands are too wide to support the claim
  you want to make, if budget allows; otherwise weaken the claim.
\end{enumerate}
The supplement is independent and can run in parallel on separate GPUs once your evaluation
harness exists.

% =====================================================================
\section{Milestones and sanity checks}\label{sec:proj5:milestones}
\begin{enumerate}
\item The environment checker passes; \texttt{uv run pytest} collects all tests.
\item Prompting evaluation runs end to end on 100 examples; manual inspection agrees with the
  grader on at least 9 of 10 sampled outputs.
\item Each component passes its tests \emph{and} your invariants (\cref{sec:proj5:grpo}); the
  1-vs-$k$ micro-batch gradient check passes.
\item The trainer's log-probabilities for a vLLM rollout match vLLM's own to within about
  $10^{-2}$ per token right after a sync; a second sync changes them after one optimizer step.
\item A 50-step run raises validation reward; entropy decreases smoothly; response lengths stay
  under the cap for most samples.
\item Four seeds of standard GRPO reach the handout's target accuracy; their spread is recorded.
\item Each sweep's plots show bands and individual runs; every claim in the write-up cites a plot.
\end{enumerate}

% =====================================================================
\section{Doing this outside Stanford}\label{sec:proj5:outside}

You need two GPUs with bf16 support. One trains a 1B-class model with AdamW (about 24\,GB plus
activations); the other serves it with vLLM and a KV pool. Two 40--80\,GB data-centre GPUs are
comfortable; two 24\,GB cards work with small micro-batches and a reduced
\code{gpu\_memory\_utilization}. With a single GPU, colocate: run vLLM with a small memory
fraction (or its sleep/offload mode, if your version supports it) and accept slower steps, or
replace vLLM with this book's engine (\cref{ch:serving}). The cost models in
\cref{sec:rlsys:placement} tell you what colocation costs. Rent cloud GPUs for the sweeps only, after
the 50-step correctness run passes locally. If even that is out of reach, reproduce every
algorithmic comparison at small scale: Lab~24's toy task and Lab~27/28's addition task run on a
CPU and show the same qualitative effects for baselines, normalisers and reuse. Then spend
GPU time only on the headline GRPO run.

% =====================================================================
\section{Summary}
\begin{itemize}
\item The project is RLVR on a small, strong base model: prompting baselines, GRPO from first
  principles, normaliser and aggregation variants, off-policy reuse, and one estimator of your own.
\item Setup is \texttt{uv} plus a pinned vLLM, two GPUs connected by an NCCL weight-sync group,
  and thin adapters into your own tested functions.
\item The mathematics is \cref{ch:pg}; the engineering invariants are the mask alignment, fp32
  log-probs, global normalisers under accumulation, and weight-sync checks.
\item Good RL science here means correctness runs before sweeps, four seeds with bands,
  one change at a time, and an honest account of tuning asymmetries.
\item The supplement walks the classic pipeline: evaluate, SFT with packing, red-team, DPO,
  evaluate again.
\end{itemize}

% =====================================================================
\section{Exercises}\label{sec:proj5:exercises}

\begin{exercise}[budgeting]
Your measured step time for standard GRPO is 40\,s on two H100s, of which 26\,s is generation.
Estimate the GPU-hours for the four-variant on-policy sweep at 200 steps and 4 seeds, and the
saving if zero-advantage pruning removes 30\% of trainer work.
\end{exercise}

\begin{exercise}[off-policy timeline]
Redraw \cref{fig:proj5:timeline} for a 32$\times$ off-policy schedule in which each optimizer
step on 1/32 of the batch takes 0.5\,s and generation takes 26\,s. What fraction of wall-clock
time is GPU~1 idle, and how many optimizer steps happen per hour compared with the on-policy run?
\end{exercise}

\begin{exercise}[reward noise]
Prompting evaluation shows that 4\% of outputs are correct but unparsed. If the grader's misses
are independent of the response, how do they bias the group-mean advantage of a group with true
success rate 0.5? What if misses are concentrated in long responses?
\end{exercise}

\begin{exercise}[single-GPU colocation]
Using the placement cost model of \cref{ch:rlsys}, estimate the slowdown of a colocated
single-GPU run compared with the two-GPU layout if weight reload and KV-pool reallocation cost
5\,s per step.
\end{exercise}

\begin{exercise}[DPO memory]
Lay out the memory of DPO training for an 8B policy with RMSprop on GPU~0 and the reference on
GPU~1, at sequence length 512 and micro-batch 1 pair. Which term would you cut first if the
policy GPU had 80\,GB?
\end{exercise}
